% Table 4 — durability is a fault-reach boundary, not a storage checklist.
% Previously a TikZ grid of colored cells with no legend; the two colors (pd
% focus fill / pd caution fill) carried no information the cell text did not
% already state, so a reader could only guess at what "durable" vs. "may
% lose last commits" being tinted differently was supposed to mean. The cell
% text is the whole claim; color was pure, unexplained redundancy. Fixed as
% a plain table -- the house pattern this chapter already uses for other
% small cross-tabs (Table~\ref{tab:swk-seven-organs}, etc.).
\begin{table}[H]
\centering
\caption{Durability by fault class under the two synchronization levels. The
reference implementation ships \texttt{synchronous=NORMAL}: a process crash
never loses a returned write, but an OS crash or power loss can roll back the
most recent commits. \texttt{synchronous=FULL} buys the power-loss cell too,
at the cost of an fsync on every commit.}
\label{tab:swk-durability-faultclass}
\small
\renewcommand{\arraystretch}{1.35}
\begin{tabularx}{\textwidth}{@{}>{\raggedright\arraybackslash}p{0.32\textwidth} >{\centering\arraybackslash}X >{\centering\arraybackslash}X >{\centering\arraybackslash}X@{}}
\toprule
 & \textbf{process crash} & \textbf{OS crash} & \textbf{power loss} \\
\midrule
\texttt{synchronous=NORMAL} & durable & may lose last commits & may lose last commits \\
\texttt{synchronous=FULL} & durable & durable & durable \\
\bottomrule
\end{tabularx}

\smallskip
\noindent\textit{The default (\texttt{NORMAL}) trades the power-loss cell for one fsync per commit.}
\end{table}
